\section{Sistemas de recomendación vs.\ sistemas generativos: sin intermediarios que se queden con la diferencia}

Esta sección entra en la teoría de plataformas más central de todo el episodio. La esencia de un sistema de recomendación es la asignación de inventario: la plataforma empareja el contenido existente con las preferencias de los usuarios y controla el ordenamiento, la distribución y la fijación del precio del tráfico. La esencia de un sistema generativo es la producción bajo demanda: la necesidad del usuario entra directamente en el extremo de producción, que genera el contenido y lo entrega directamente al usuario; la etapa de asignación queda interiorizada en el sistema de producción.

\begin{figure}[H]
\centering
\includegraphics[width=0.94\textwidth]{recommendation-vs-generation.png}
\caption{Sistemas de recomendación vs.\ sistemas generativos: la diferencia de fondo entre la asignación de inventario y la producción bajo demanda. Diagrama conceptual propio, elaborado a partir de la conversación de 02:38:11--02:50:34.}
\end{figure}

\begin{knowledgebox}{Lectura de la figura: el intermediario es el «derecho de asignación del inventario»}
El sistema de recomendación empareja contenido a partir del inventario existente, y la plataforma controla el derecho de distribución; el sistema generativo produce al instante según la necesidad del usuario, y la asignación queda interiorizada en la producción. La expresión «sin intermediarios que se queden con la diferencia» no significa que no haya plataforma, sino que el derecho de distribución de las plataformas tradicionales deja de ser la única forma de convertir el valor en ingresos.
\end{knowledgebox}

\subsection{Por qué el sistema de recomendación es un intermediario}

Este apartado explica primero el sentido mecánico de la palabra «intermediario». No es una crítica moral: significa que la plataforma controla la asignación del inventario, el ordenamiento y el poder de comercialización, y por eso puede extraer valor entre productores y consumidores.

El poder de las plataformas de internet proviene de una relación bilateral: los creadores producen contenido, los consumidores lo consumen y la plataforma se encarga del emparejamiento y extrae valor. Los motores de búsqueda, los motores de recomendación y las plataformas de comercio electrónico actúan, en distinto grado, como intermediarios. Los creadores invierten recursos en producir contenido, pero el mismo contenido publicado desde cuentas distintas puede recibir un tráfico completamente diferente; esto muestra que el mecanismo de asignación no depende solo del contenido en sí, sino también del peso que la plataforma asigna, del historial de la cuenta, de los objetivos del algoritmo y de las reglas comerciales.

\begin{warningbox}{No hay que entender «la plataforma es un intermediario» como que la plataforma no aporta valor}
Las plataformas sí crean valor en descubrimiento, emparejamiento, prevención de fraude, pagos, comunidad y distribución a escala. El problema es que los sistemas generativos reducen la escasez relativa de la «asignación de inventario» y obligan a las plataformas a desplazarse hacia el sistema de producción, las redes de confianza u otras capas de valor.
\end{warningbox}

\subsection{Cómo interioriza la asignación el sistema generativo}

El sistema generativo no busca en el inventario el contenido más cercano a la necesidad del usuario, sino que produce el contenido directamente a partir de la intención del usuario. Sigue habiendo asignación, pero ocurre dentro del proceso de generación: la necesidad y las preferencias del usuario, el context, la recipe y las capacidades del sistema determinan en conjunto la salida. Así, la plataforma deja de ser la única que controla si un contenido se ve o no, y la intención del usuario se convierte en una entrada más directa de la producción.

\begin{figure}[H]
\centering
\includegraphics[width=0.96\textwidth]{power-shift-to-consumer.png}
\caption{El poder se desplaza hacia el lado del consumidor: la cadena de los bienes de información pasa de la asignación por la plataforma a estar impulsada por la intención del usuario. Diagrama conceptual propio, elaborado a partir de la conversación de 02:25:49--02:38:11.}
\end{figure}

\begin{knowledgebox}{Lectura de la figura: el consumidor no es solo el extremo receptor}
En la cadena tradicional, el productor produce, la plataforma distribuye y el consumidor recibe. El sistema generativo hace que la intención del consumidor entre en el extremo de producción: el usuario no solo consume contenido de forma pasiva, sino que, al consumirlo, participa en su producción.
\end{knowledgebox}

\subsection{Resumen de la sección}

La diferencia de fondo entre los sistemas de recomendación y los sistemas generativos es la diferencia entre la asignación de inventario y la producción bajo demanda. Las plataformas no desaparecerán, pero la forma en que las plataformas tradicionales extraen valor mediante el derecho de distribución se verá cuestionada.
